\begin{figure}[!t]
  \centering
  \makebox[\linewidth][c]{%
    \resizebox{1.08\linewidth}{!}{%
      \begin{tikzpicture}[
        x=1cm,
        y=1cm,
        phasebox/.style={draw=black, line width=0.8pt, fill=white,
          rounded corners=1pt, minimum width=72mm, minimum height=31mm,
          text width=66mm, inner sep=4pt, align=center, font=\large},
        phaseviolet/.style={phasebox},
        phaseorange/.style={phasebox},
        qrombox/.style={draw=black, line width=0.8pt, fill=orange!12,
          rounded corners=1pt, minimum width=38mm, minimum height=13mm,
          inner sep=3pt, align=center, font=\large},
        wire/.style={line width=1pt},
        statebox/.style={draw=black!55, line width=0.7pt, fill=white,
          rounded corners=1pt, inner sep=5pt, align=center, font=\large},
      ]
        \path[use as bounding box] (-1.5,0) rectangle (27.4,16.0);

        % Inputs and transient addend materialization.
        \node[font=\Large, anchor=west] at (0.25,15.45)
          {Input: $R=(x_R,y_R)$, selected addend $A=(x_A,y_A)$};
        \node[qrombox] (qAin) at (4.2,13.65)
          {load/unload $A$\\$A=(x_A,y_A)$};
        \node[qrombox] (q3x) at (21.6,13.65)
          {load/unload $3x_A$};

        % First row: algorithm lines 1--3.
        \node[phasebox] (s1) at (4.2,10.65)
          {\textbf{1. Coordinate differences}\\[-0.5mm]
           \textit{\Cref{alg:point-addition} line 1}\\[1mm]
           $X\gets x_R-x_A=d_x$\\
           $Y\gets y_R-y_A=d_y$};
        \node[phaseviolet] (s2) at (12.9,10.65)
          {\textbf{2. Dialog-GCD division}\\[-0.5mm]
           \textit{\Cref{alg:point-addition} line 2}\\[1mm]
           $Y\gets d_y d_x^{-1}=\lambda\pmod p$\\
           $X=d_x$; clean GCD workspace};
        \node[phasebox] (s3) at (21.6,10.65)
          {\textbf{3. Prepare the $X$ workspace}\\[-0.5mm]
           \textit{\Cref{alg:point-addition} line 3}\\[1mm]
           $X\gets d_x+3x_A=x_R+2x_A$\\
           $Y=\lambda$; then uncompute $3x_A$};
        \draw[->, dashed, line width=0.8pt] (qAin.south) -- (s1.north);
        \draw[->, dashed, line width=0.8pt] (q3x.south) -- (s3.north);

        % Second row: algorithm lines 4--7.
        \node[phaseorange] (s4) at (4.2,5.55)
          {\textbf{4. Specialized modular square}\\[-0.5mm]
           \textit{\Cref{alg:point-addition} line 4}\\[1mm]
           $X\gets X-\lambda^2=x_A-x_{R'}$\\
           $Y=\lambda$};
        \node[phaseviolet] (s5) at (12.9,5.55)
          {\textbf{5. Forward multiplication}\\[-0.5mm]
           \textit{\Cref{alg:point-addition} line 5}\\[1mm]
           $Y\gets\lambda X=\lambda(x_A-x_{R'})$\\
           $X=x_A-x_{R'}$; erase transcript};
        \node[phasebox] (s6) at (21.6,5.55)
          {\textbf{6. Recover the output}\\[-0.5mm]
           \textit{\Cref{alg:point-addition} lines 6--7}\\[1mm]
           $Y\gets Y-y_A=y_{R'}$\\
           $X\gets x_A-X=x_{R'}$\\[1mm]
           Output: $(X,Y)=(x_{R'},y_{R'})$};
        % The recycled X and Y registers pass through every phase.  The two
        % wires wrap from the end of the first row to the beginning of the
        % second row; drawing them on the background layer hides the portions
        % that pass through the filled phase boxes.
        \begin{scope}[on background layer]
          \draw[wire]
            (-0.05,11.15) -- (26.00,11.15) -- (26.00,8.55)
            -- (-0.05,8.55) -- (-0.05,6.05) -- (26.15,6.05);
          \draw[wire]
            (-0.05,10.15) -- (26.40,10.15) -- (26.40,8.10)
            -- (-0.45,8.10) -- (-0.45,5.05) -- (26.15,5.05);
        \end{scope}

        % Input and output register labels.
        \node[font=\Large, anchor=east] at (-0.15,11.15)
          {$X:\ket{x_R}$};
        \node[font=\Large, anchor=east] at (-0.15,10.15)
          {$Y:\ket{y_R}$};
        \node[font=\Large, anchor=west] at (26.25,6.05)
          {$X=\ket{x_{R'}}$};
        \node[font=\Large, anchor=west] at (26.25,5.05)
          {$Y=\ket{y_{R'}}$};

        % The final recovery temporarily rematerializes A and then clears it.
        \node[qrombox] (qAout) at (21.6,2.75)
          {load/unload $A$\\recover $R'$ and return $A$ to zero};
        \draw[->, dashed, line width=0.8pt] (qAout.north) -- (s6.south);

        % Compact summary of the reusable machinery.
        \node[statebox, text width=230mm, minimum height=13mm]
          at (12.9,0.95)
          {Stages 2 and 5 reuse the dialog-GCD machinery: division produces
           $\lambda=d_y/d_x$, while reverse use of the same reversible map
           multiplies by $x_A-x_{R'}$ and clears the transcript.};
      \end{tikzpicture}%
    }%
  }
  \caption{High-level architecture of the best $Q\times T$-scoring
  \texttt{8e9c9a2} mixed point-addition circuit. The six phases implement the
  in-place map $R\mapsto R'=R+A$ while recycling the $X$ and $Y$ coordinate
  registers. Under the challenge evaluator, the circuit uses $Q=1{,}151$
  logical qubits and $T=1{,}299{,}453$ average executed Toffoli gates. The
  lookup/unlookup boxes indicate where a fully windowed Shor construction
  would temporarily materialize the selected addend; in the challenge
  benchmark, $(x_A,y_A)$ are provided as classical inputs.}
  \label{fig:best-qxt-circuit-pipeline}
\end{figure}
